\documentclass[11pt]{article}

\usepackage[margin=1in]{geometry}
\usepackage{amsmath,amssymb}
\usepackage{booktabs}
\usepackage{tabularx}
\usepackage{array}
\usepackage{enumitem}
\usepackage[T1]{fontenc}
\usepackage{hyperref}

\newcommand{\code}[1]{\texttt{#1}}
\newcommand{\e}{\mathbf{e}}
\newcommand{\rt}{\tilde\rho}
\newcommand{\rb}{\bar\rho}
\newcommand{\Wh}{\hat W}
\newcolumntype{L}{>{\raggedright\arraybackslash}X}
\emergencystretch=3em

\title{The flocking model in \code{src\_deankow/flocking}:\\
formulation, stability and sensing kernels}
\author{FHDeX}
\date{}

\begin{document}
\maketitle

\noindent This note describes the model solved by the codes in
\code{src\_deankow/flocking}:
\begin{itemize}[nosep]
  \item active Brownian particles moving at a speed that can depend on the
        local smoothed density;
  \item the Dean--Kawasaki SPDE for their empirical density;
  \item the linear stability of the uniform state;
  \item the choice of sensing kernel.
\end{itemize}
\code{README.md} in the same directory covers the code layout, inputs and
numerical checks in detail. The density is written $\phi$ in the code; it is
$\rho$ in \code{src\_deankow/flocking.pdf}.

\section{Particle model}

$N$ particles move in the periodic box $[0,L_x)\times[0,L_y)$. Each particle
has a position $\mathbf{X}_i$ and a heading $\theta_i\in[0,2\pi)$:
\begin{equation}
  d\mathbf{X}_i = v(\mathbf{X}_i)\,\e(\theta_i)\,dt, \qquad
  \e(\theta) = (\cos\theta, \sin\theta), \qquad
  d\theta_i = \sqrt{2D}\,dW_i ,
\end{equation}
with independent standard Brownian motions $W_i$ and rotational diffusivity
$D$ (\code{diff\_coeff}). The speed is
\begin{equation}
  v(\mathbf{x}) = v_0\, s\bigl(\rt(\mathbf{x})/\rb\bigr), \qquad
  \rt(\mathbf{x}) = \frac{1}{N}\sum_j W(\mathbf{x}-\mathbf{X}_j), \qquad
  \int W\,dA = 1 ,
\end{equation}
with self-propulsion speed $v_0$ (\code{flock.speed}), a dimensionless speed
function $s(u)\le 1$, a sensing kernel $W$ of range $R$, and
$\rb = 1/(L_xL_y)$ the mean of $\rt$.
\begin{itemize}
  \item With $s\equiv 1$ the particles are independent active Brownian
        particles.
  \item With $s$ decreasing they slow down in crowded regions, which can lead
        to motility-induced phase separation (MIPS).
\end{itemize}
Positions are moved deterministically given the headings; all the noise is in
$\theta$. The coupling is mean-field: each particle contributes $W/N$ to
$\rt$.

\paragraph{Scales.}
\begin{itemize}
  \item Persistence time $1/D$: a particle loses memory of its heading on this
        time.
  \item Persistence length $\ell_p = v/D$: the distance travelled before the
        heading decorrelates.
  \item Long-time diffusivity $D_{\mathrm{eff}} = v^2/(2D)$ in 2D: on scales
        $\gg\ell_p$, a particle at constant speed $v$ random-walks with this
        diffusivity.
\end{itemize}
The dimensionless groups are $L/\ell_p$, $R/\ell_p$, the speed-function
parameter $\lambda$, and $N$, which sets the size of fluctuations.

\section{Phase-space density}

Let
\begin{equation}
  \phi(\mathbf{x},\theta,t) = \frac1N\sum_i
     \delta\bigl(\mathbf{x}-\mathbf{X}_i(t)\bigr)\,\delta\bigl(\theta-\theta_i(t)\bigr),
  \qquad \iint\phi\,d\theta\,d\mathbf{x} = 1 ,
\end{equation}
with spatial density $\rho = \int\phi\,d\theta$, polarization
$\mathbf{p} = \int\e(\theta)\,\phi\,d\theta$, and smoothed density
$\rt = W\star\rho$. In the code, $\phi$ is a cell-averaged field on a 3D grid
$(x,y,\theta)$, with $\theta$ the $z$ direction.

\subsection{Mean-field (kinetic) equation}
As $N\to\infty$, $\phi$ obeys the nonlinear Fokker--Planck equation
\begin{equation}
  \partial_t\phi = D\,\partial_{\theta\theta}\phi
     - \nabla_{\mathbf{x}}\cdot\bigl(v(\rt)\,\e(\theta)\,\phi\bigr),
  \qquad v(\rt) = v_0\,s(\rt/\rb).
  \label{eq:kinetic}
\end{equation}

\subsection{Dean--Kawasaki SPDE}
For finite $N$, Itô's formula applied to the empirical measure gives,
formally (Dean 1996; Müller, von Renesse and Zimmer 2025),
\begin{equation}
  d\phi = D\,\partial_{\theta\theta}\phi\,dt
    - \nabla_{\mathbf{x}}\cdot\bigl(v(\rt)\,\e(\theta)\,\phi\bigr)\,dt
    + N^{-1/2}\,\partial_\theta\bigl(\sqrt{2D\phi}\,dB\bigr),
  \label{eq:dk}
\end{equation}
where $B$ is space--angle--time white noise. For a test function
$f(\mathbf{x},\theta)$,
\begin{equation}
  d\langle\phi,f\rangle = \bigl\langle\phi,\ v\,\e\cdot\nabla_{\mathbf{x}}f
      + D\,\partial_{\theta\theta}f\bigr\rangle\,dt
    + \frac1N\sum_i\sqrt{2D}\,\partial_\theta f(\mathbf{X}_i,\theta_i)\,dW_i ,
\end{equation}
and the martingale term has quadratic variation
$(2D/N)\langle\phi,(\partial_\theta f)^2\rangle\,dt$. That is exactly the
conservative noise in \eqref{eq:dk}.
\begin{itemize}
  \item There is no noise in $\mathbf{x}$, because the spatial motion is
        deterministic given the headings.
  \item The interaction adds no noise: it acts only through $v(\rt)$, itself
        a functional of $\phi$.
\end{itemize}
The noise conserves mass at each spatial point: it moves probability in
$\theta$ only.

\subsection{Moment equations and the diffusive limit}
The $\theta$-moments of \eqref{eq:kinetic} are
\begin{equation}
  \partial_t\rho = -\nabla\cdot(v\,\mathbf{p}), \qquad
  \partial_t\mathbf{p} = -D\,\mathbf{p}
     - \nabla\cdot\Bigl(v\int\e\,\e^{\mathsf T}\phi\,d\theta\Bigr).
\end{equation}
On length scales $\gg\ell_p$ and time scales $\gg 1/D$, the polarization
relaxes quickly and $\int\e\e^{\mathsf T}\phi\,d\theta\approx(\rho/2)\,I$.
Then $\mathbf{p}\approx-\nabla(v\rho)/(2D)$ and
\begin{equation}
  \partial_t\rho = \nabla\cdot\Bigl(\frac{v}{2D}\,\nabla(v\rho)\Bigr).
\end{equation}
For a local speed $v(\rho)$ this is
\begin{equation}
  \partial_t\rho = \nabla\cdot\Bigl(D_{\mathrm{eff}}(\rho)
     \Bigl[1+\frac{d\ln v}{d\ln\rho}\Bigr]\nabla\rho\Bigr),
  \qquad D_{\mathrm{eff}}(\rho) = \frac{v(\rho)^2}{2D}.
\end{equation}
The collective diffusivity turns negative when $d\ln v/d\ln\rho<-1$. That is
the spinodal condition for MIPS (Tailleur and Cates 2008; Cates and Tailleur
2015). Particles accumulate where they are slow, and if they slow down fast
enough with density, the accumulation runs away.

\section{Linear stability of the uniform state}

The uniform state $\phi_0 = \rb/(2\pi)$ is stationary for any $s$ and $W$.
With $u = \rt/\rb$, define at the mean density $u=1$
\begin{equation}
  v_1 = v_0\,s(1), \qquad
  g = \left.\frac{d\ln s}{d\ln u}\right|_{u=1} = \frac{s'(1)}{s(1)} .
\end{equation}
Perturbing $\phi=\phi_0+\delta\phi$ gives $\delta\rt = W\star\delta\rho$ and
$\delta v = v_1\,g\,\delta\rt/\rb$.

\subsection{Exact (kinetic) dispersion relation}
Take a mode $e^{i\mathbf{k}\cdot\mathbf{x}}$ with $\mathbf{k}$ along $x$, and
write $\delta\phi = \sum_m a_m e^{im\theta}$, so that $\delta\rho = 2\pi a_0$.
The linearized equation is
\begin{equation}
  \partial_t\delta\phi = D\,\partial_{\theta\theta}\delta\phi
     - ik\,v_1\cos\theta\,\bigl[\delta\phi + g\,\Wh(k)\,a_0\bigr],
\end{equation}
with $\Wh$ the 2D Fourier transform of $W$. In harmonics,
\begin{align}
  \partial_t a_0 &= -\tfrac{ikv_1}{2}\,(a_1+a_{-1}),\\
  \partial_t a_{\pm1} &= -D\,a_{\pm1}
      - \tfrac{ikv_1}{2}\bigl[(1+g\Wh(k))\,a_0 + a_{\pm2}\bigr],\\
  \partial_t a_m &= -Dm^2 a_m - \tfrac{ikv_1}{2}\,(a_{m-1}+a_{m+1}),
      \qquad |m|\ge2 .
\end{align}
The interaction enters only through the $m=0\to\pm1$ coupling, multiplied by
$1+g\Wh(k)$. The growth rate $\sigma(k)$ is the eigenvalue of this
tridiagonal system with the largest real part; it can also be written as a
continued fraction. The SPDE resolves $\theta$ with $n_\theta$ cells, so the
code's exact rate is the leading eigenvalue of the same operator with
$\partial_{\theta\theta}$ replaced by the $n_\theta$-point second difference.

\subsection{Diffusive limit}
For $k\ell_1\ll1$, with $\ell_1 = v_1/D$, the harmonics with $|m|\ge2$ are
negligible and $a_{\pm1}\approx-(ikv_1/2D)(1+g\Wh)\,a_0$. That gives
\begin{equation}
  \boxed{\ \sigma(k) = -D_{\mathrm{eff}}\,k^2\,\bigl(1+g\,\Wh(k)\bigr),
  \qquad D_{\mathrm{eff}} = \frac{v_1^2}{2D}\ }
\end{equation}
which is the Fourier form of the diffusive limit above with the nonlocal
speed. The uniform state is unstable at $k$ when $1+g\Wh(k)<0$.
\begin{itemize}
  \item \textbf{Onset.} $\Wh(0)=1$, so long waves go unstable first. The
        onset is $g=-1$: $d\ln v/d\ln\rho=-1$ at the mean density.
  \item \textbf{Unstable band.} For a decreasing $\Wh$ (Gaussian) the band is
        $0<k<k_c$ with $\Wh(k_c)=-1/g$. The sensing range sets the
        short-wavelength cutoff, $k_c\sim1/R$.
  \item \textbf{Fastest mode.} $\sigma$ is largest where
        $k^2(-1-g\Wh(k))$ peaks. That sets the initial domain spacing, which
        later coarsens in the nonlinear regime.
  \item \textbf{Mass conservation.} $\sigma\to0$ as $k\to0$, and the $k=0$
        mode (total mass) is neutral.
  \item \textbf{Kinetic corrections.} At $k\ell_1\approx0.5$--$1$ the
        diffusive formula overestimates the growth, because finite
        persistence lags the response. The startup printout reports the
        diffusive rates and $k\ell_p$ for the first four box modes.
\end{itemize}

\subsection{Speed functions}
\begin{center}
\begin{tabular}{@{}lllll@{}}
\toprule
\code{speed\_type} & $s(u)$ & $s(1)$ & $g$ & long-wave instability \\
\midrule
0 & $1$ & $1$ & $0$ & never (independent particles) \\
1 & $e^{-\lambda u}$ & $e^{-\lambda}$ & $-\lambda$ & $\lambda>1$ \\
2 & $\max(0,1-\lambda u)$, $\lambda<1$ & $1-\lambda$ & $-\lambda/(1-\lambda)$ & $\lambda>1/2$ \\
\bottomrule
\end{tabular}
\end{center}
\begin{itemize}
  \item \textbf{Exponential:} stays positive at any density. Stronger
        $\lambda$ also lowers the speed at the mean density,
        $v_1 = v_0e^{-\lambda}$, which reduces $D_{\mathrm{eff}}$ and slows
        the dynamics.
  \item \textbf{Linear:} stops particles completely at $u=1/\lambda$. It
        needs $\lambda<1$, so that the mean-density speed is positive.
\end{itemize}

\subsection{Checks of the code against the theory}
The deterministic SPDE (\code{dorand = 0}, centred advection) was run with
$v_0=4$, $D=10$, a Gaussian kernel with $R=0.03$, and mode $k=2\pi$. Its
measured growth rates compare as follows:
\begin{center}
\begin{tabular}{@{}lccc@{}}
\toprule
$\lambda$ & SPDE & kinetic eigenvalue (32-point $\theta$) & diffusive limit \\
\midrule
0.8 (stable) & $-1.448$ & $-1.462$ & $-1.365$ \\
1.5 (unstable) & $+0.684$ & $+0.686$ & $+0.745$ \\
\bottomrule
\end{tabular}
\end{center}
$k\ell_1$ is 1.1 and 0.56 for these two runs. The SPDE matches the kinetic
theory to 1\%; the diffusive formula is 6--9\% off at these $k\ell_1$.

\subsection{The non-interacting case}
With $s\equiv1$ the equation is linear, and every perturbation decays. The
$\theta$ harmonics damp at rates $Dm^2$. Spatial structure relaxes
diffusively with $D_{\mathrm{eff}}$ on scales $\gg\ell_p$, and ballistically
below $\ell_p$.

The fluctuations are Poisson. The stationary covariance of the linear SPDE is
\begin{equation}
  \langle\delta\phi(\mathbf{x},\theta)\,\delta\phi(\mathbf{x}',\theta')\rangle
  = \frac{\phi_0}{N}\,\delta(\mathbf{x}-\mathbf{x}')\,\delta(\theta-\theta') :
\end{equation}
the $\theta$ diffusion and $\theta$ noise balance, and the advection
$-\e\cdot\nabla_{\mathbf{x}}$ is skew-adjoint, so it neither creates nor
destroys variance.

The discretization must keep this property. Centred advection does, and the
code reproduces Poisson statistics (variance ratio 1.00). Upwind and limited
(MUSCL) fluxes add numerical diffusion without matching noise. They damp the
spatial density variance to about 5\% of the correct value. This is why the
stochastic inputs use centred advection with the stochastic RK3 integrator.

\section{Choice of sensing kernel}

$W$ is radially symmetric, periodic, and normalized to $\int W\,dA=1$; by the
dispersion relation, stability depends on $W$ only through $\Wh(k)$. Four
kernels are available (\code{flock.kernel\_type}, radius
$R$ = \code{flock.kernel\_R}):
\begin{center}
\begin{tabularx}{\textwidth}{@{}lllL@{}}
\toprule
type & $W(r)\propto$ (normalization) & $\Wh(k)$ & properties \\
\midrule
0 Gaussian & $e^{-r^2/2R^2}$ \ $(1/2\pi R^2)$ & $e^{-k^2R^2/2}$ & $C^\infty$; $\Wh>0$, monotone; cut off at $4R$ for the pair sum \\
1 top-hat & $1$, $r<R$ \ $(1/\pi R^2)$ & $2J_1(kR)/(kR)$ & discontinuous; ``neighbours within $R$''; lobe down to about $-0.13$ \\
2 parabolic & $1-r^2/R^2$, $r<R$ \ $(2/\pi R^2)$ & $8J_2(kR)/(kR)^2$ & continuous, kink at $R$; same shape as \code{hk} \\
3 bump & $(1-r^2/R^2)^2$, $r<R$ \ $(3/\pi R^2)$ & $48J_3(kR)/(kR)^3$ & $C^1$, compact; small lobes \\
\bottomrule
\end{tabularx}
\end{center}
\begin{itemize}
  \item \textbf{Onset is kernel-independent.} It is set by $g$ at $k\to0$,
        because $\Wh(0)=1$ for every kernel.
  \item \textbf{The unstable band and fastest mode depend on the kernel.}
        They follow from $1+g\Wh(k)<0$. Larger $R$ narrows the band toward
        long waves and gives larger initial domains.
  \item \textbf{Negative lobes.} For a decreasing $s$ ($g<0$), a negative
        $\Wh$ makes $1+g\Wh>1$, so the lobes of the compact kernels are
        stabilizing. They cannot create a finite-wavelength instability, but
        would matter for an increasing speed function ($g>0$). The Gaussian
        gives the cleanest long-wave instability.
  \item \textbf{Grid resolution.} On the grid, $W$ is sampled at the cell
        offsets and normalized discretely, $\sum W\,\Delta x\,\Delta y=1$, so
        that the mean of $\rt$ is exactly $\rb$. A $\cos kx$ mode is then
        damped by the discrete transform, which matches $\Wh$ to
        $O((\Delta x/R)^2)$ for smooth kernels. Measured: 0.951975 against
        0.951850 for the Gaussian at $R=3.2\,\Delta x$. Keep $R$ at about
        $4\,\Delta x$ or more; the top-hat, with its edge, needs more.
  \item \textbf{Particles: particle-mesh vs pair sum.} Particle-mesh uses a
        cloud-in-cell deposit, the discrete $W$, and cloud-in-cell
        interpolation, which adds about one cell of smoothing. Against the
        exact pair sum (bump kernel, same particles), the rms difference is
        1.4\% at $R=2.6\,\Delta x$, 0.15\% at $5.1\,\Delta x$ and 0.017\% at
        $10.2\,\Delta x$: about $(\Delta x/R)^3$.
  \item \textbf{Cost of the pair sum.} It grows like $N\pi R^2/A$ neighbours
        per particle, so compact kernels keep it bounded. The Gaussian's $4R$
        cutoff costs about $16\times$ the neighbours of a compact kernel of the
        same $R$.
\end{itemize}
\paragraph{Recommendations.}
\begin{itemize}
  \item SPDE and particle-mesh runs: the Gaussian, for its positive, monotone
        $\Wh$ and its smoothness.
  \item Pair-sum particle runs: the bump.
  \item Comparison with ``neighbours within $R$'' models in the literature:
        the top-hat.
\end{itemize}

\section{Numerical formulation}

\paragraph{SPDE} (3D build; \code{AmrCoreFlock}, \code{mykernel.H}). A
finite-volume scheme on $(x,y,\theta)$ cells, with $\theta_k$ the cell-centre
heading. The fluxes are:
\begin{itemize}
  \item $x$, $y$ faces: advection $\phi\,v_{\mathrm{face}}(\cos\theta_k,\sin\theta_k)$,
        with $v_{\mathrm{face}} = v_0\,s\bigl(\tfrac12(\rt_L+\rt_R)/\rb\bigr)$.
        The flux is centred (\code{adv\_order = 0}), first-order upwind (1),
        or limited MUSCL (2).
  \item $\theta$ faces: $-D\,\partial_\theta\phi$ (centred), plus the noise
        $\sqrt{2D\phi_{\mathrm{face}}/(N\,\Delta V\,\Delta t)}\,Z$ with
        $Z\sim N(0,1)$.
\end{itemize}
$\rt$ is a periodic FFT convolution of $\phi$ with $W$, replicated over
$\theta$ (\code{DensityConv}): $\phi$ is integrated over $\theta$ and convolved on
the $x$--$y$ plane by a 2D FFT. It is recomputed for each stage of the time
integrator. The options are Euler--Maruyama (\code{time\_integrator = 0}),
Heun (1), or the stochastic SSP-RK3 of Delong, Griffith, Vanden-Eijnden and
Donev (2013) (2). RK3 is stable for centred advection and third order for the
deterministic part.

The time step is $\Delta t = \mathrm{cfl}/\bigl(v_0(1/\Delta x+1/\Delta y)
+ 2D/\Delta\theta^2\bigr)$, the positivity limit of upwind with forward
Euler; $v_0$ bounds the speed because $s\le1$. The $\theta$-diffusion term
grows like $n_\theta^2$, so with fine $\theta$ grids an implicit $\theta$
diffusion would be the next improvement.

\paragraph{Particles} (2D build; \code{FlockPC}).
\begin{itemize}
  \item Euler--Maruyama: the position is advanced with the speed and heading
        at the start of the step, and $\theta \mathrel{+}= \sqrt{2D\Delta t}\,Z$.
  \item $\rt$ comes from particle-mesh (\code{density\_method = 0}) or the
        pair sum (1).
\end{itemize}
The particle histogram over $\theta$ bins on the SPDE grid is compared cell
for cell with the SPDE by \code{python/flocking\_compare.py}.

\section{Parameter regimes}

MIPS requires:
\begin{itemize}
  \item $L\gg\ell_1 = v_1/D$, so the box holds many persistence lengths;
  \item $g<-1$, for example $\lambda>1$ for the exponential;
  \item $R$ small enough that the unstable band contains several box modes,
        but at least about $4\,\Delta x$;
  \item $N$ large enough that a kernel area holds many particles,
        $N\pi R^2/A\gg1$.
\end{itemize}
The example \code{inputs\_mips} uses $v_0=4$, $D=10$, $\lambda=1.5$, a
Gaussian kernel with $R=0.03$, a $128^2\times32$ grid and $N=10^6$. That
gives $v_1\approx0.89$, $\ell_1\approx0.09$, and box modes 1--4 unstable.
Starting from uniform, both codes form dense, slow domains about a third of
the box apart. By $t\approx3$ about 26\% of the area has $\rt>1.5\rb$ and
about 55\% has $\rt<0.5\rb$.

\section{Outlook: alignment}

Flocking proper adds a turning rate to the heading equation,
\begin{equation}
  d\theta_i = \omega(\mathbf{X}_i,\theta_i)\,dt + \sqrt{2D}\,dW_i, \qquad
  \omega = -K\,\operatorname{Im}\bigl(e^{-i\theta}(W\star P)(\mathbf{x})\bigr),
  \qquad P = \int e^{i\theta}\phi\,d\theta ,
\end{equation}
which is a Vicsek/Kuramoto alignment. In the SPDE this is a $\theta$-advection
term $-\partial_\theta(\omega\phi)$. The kernel already has a slot for it in
the $\theta$ flux, and $\omega$ needs one more FFT convolution of the first
$\theta$-moment. Mean-field theory then predicts:
\begin{itemize}
  \item an isotropic-to-polar transition near $K\rb/2\approx D$, with the
        exact threshold depending on how $K$ is normalized;
  \item travelling bands near the transition;
  \item giant number fluctuations in the ordered phase.
\end{itemize}
These are a natural next test of whether the SPDE reproduces finite-$N$
effects.

\begin{thebibliography}{9}
\bibitem{dean} D.~S.~Dean, Langevin equation for the density of a system of
  interacting Langevin processes, \emph{J. Phys. A} \textbf{29}, L613 (1996).
\bibitem{tc08} J.~Tailleur and M.~E.~Cates, Statistical mechanics of
  interacting run-and-tumble bacteria, \emph{Phys. Rev. Lett.} \textbf{100},
  218103 (2008).
\bibitem{ct15} M.~E.~Cates and J.~Tailleur, Motility-induced phase separation,
  \emph{Annu. Rev. Condens. Matter Phys.} \textbf{6}, 219 (2015).
\bibitem{mrz} Müller, von Renesse and Zimmer (2025), as cited in
  \code{src\_deankow/flocking.pdf}, for the Dean--Kawasaki equation of this
  active-particle model.
\end{thebibliography}

\end{document}
